% GENERATED FILE. Edit canonical lessons, then run npm run generate:books.
% canonical-source-hash: fnv1a64:852fa73cf9c8c3e7
% canonical-lessons: BN-C324-dhadha, BN-C324-koushol, BN-C324-phad, BN-C324-bholano, BN-C324-par-paoya

\chapter{Riddles and Traps}
\label{ch:bn-riddles-and-traps}
\hlchaptermodality{324}

\begin{chapteropening}
\textbf{By the end of this chapter:} \emph{I can tell how someone was tricked, trapped, and whether they got away.}

Complete the last lesson of chapter 324: I can tell how someone was tricked, trapped, and whether they got away.
\end{chapteropening}

\section[\texorpdfstring{\bn{ধাঁধা} (dhā̃dhā)}{ধাঁধা (dhā̃dhā)}]{\bn{ধাঁধা} (dhā̃dhā) — a riddle; a puzzle}

\label{lesson:BN-C324-dhadha}



\begin{quote}
Before the new one: say the Bengali for a holy man, then the Bengali for a disguise.
\end{quote}

\subsection*{You'll want to know: \bn{ধাঁধা}}
\textbf{\bn{ধাঁধা}} — \emph{dhā̃dhā} — \textquotedblleft{}a riddle; a puzzle\textquotedblright{}.

\begin{quote}
\bn{রাক্ষস} \bn{রাজকুমারকে} \bn{একটা} \bn{ধাঁধা} \bn{জিজ্ঞাসা} \bn{করল।}
\end{quote}

\begin{quote}
\emph{rākkhôs rājkumārke ekṭā dhā̃dhā jigyāsā kôrlô.}
\end{quote}

\begin{quote}
The ogre asked the prince a riddle.
\end{quote}

\begin{cousinweb}
\textbf{\bn{ধাঁধা}} (\emph{dhā̃dhā}) is a riddle or a puzzle; its origin is not settled. \bn{ধাঁধায়} \bn{পড়া} is to be puzzled.
\end{cousinweb}

\begin{grammarlens}[title={What you've built}]
One more word for telling what happened.
\end{grammarlens}

\subsection*{Your turn}
\begin{itemize}
\raggedright
  \item \emph{Say it:} \emph{dhā̃dhā}
  \item \emph{Say it:} \emph{dhā̃dhā}, once more
  \item \emph{Say it:} \emph{rākkhôs rājkumārke ekṭā dhā̃dhā jigyāsā kôrlô.} — The ogre asked the prince a riddle.
  \item \emph{Recall:} say the Bengali for a holy man, then the Bengali for a disguise, then say \emph{dhā̃dhā} again
\end{itemize}

\begin{tcolorbox}[breakable,colback=teal!4,colframe=teal!35!black,title={Before you move on}]
What does \bn{ধাঁধা} mean? (\textquotedblleft{}A riddle; a puzzle\textquotedblright{}.) Say it once more.
\end{tcolorbox}
\section[\texorpdfstring{\bn{কৌশল} (koushôl)}{কৌশল (koushôl)}]{\bn{কৌশল} (koushôl) — a trick; a clever plan}

\label{lesson:BN-C324-koushol}



\begin{quote}
Before the new one: say the Bengali for a disguise, then the Bengali for a riddle.
\end{quote}

\subsection*{You'll want to know: \bn{কৌশল}}
\textbf{\bn{কৌশল}} — \emph{koushôl} — \textquotedblleft{}a trick; a clever plan\textquotedblright{}.

\begin{quote}
\bn{শিয়াল} \bn{একটা} \bn{কৌশল} \bn{করে} \bn{কুমিরকে} \bn{বোকা} \bn{বানাল।}
\end{quote}

\begin{quote}
\emph{shiyāl ekṭā koushôl kôre kumirke bokā bānālô.}
\end{quote}

\begin{quote}
The jackal played a trick and made a fool of the crocodile.
\end{quote}

\begin{cousinweb}
\textbf{\bn{কৌশল}} (\emph{koushôl}) is Sanskrit \emph{kauśala}, 'welfare; skilfulness, cleverness' (Monier-Williams); in Bengali it is the clever plan itself. Its \bn{ৌ} is read \emph{ou}.
\end{cousinweb}

\begin{grammarlens}[title={What you've built}]
One more word for telling what happened.
\end{grammarlens}

\subsection*{Your turn}
\begin{itemize}
\raggedright
  \item \emph{Say it:} \emph{koushôl}
  \item \emph{Say it:} \emph{koushôl}, once more
  \item \emph{Say it:} \emph{shiyāl ekṭā koushôl kôre kumirke bokā bānālô.} — The jackal played a trick and made a fool of the crocodile.
  \item \emph{Recall:} say the Bengali for a disguise, then the Bengali for a riddle, then say \emph{koushôl} again
\end{itemize}

\begin{tcolorbox}[breakable,colback=teal!4,colframe=teal!35!black,title={Before you move on}]
What does \bn{কৌশল} mean? (\textquotedblleft{}A trick; a clever plan\textquotedblright{}.) Say it once more.
\end{tcolorbox}
\section[\texorpdfstring{\bn{ফাঁদ} (phā̃d)}{ফাঁদ (phā̃d)}]{\bn{ফাঁদ} (phā̃d) — a trap}

\label{lesson:BN-C324-phad}



\begin{quote}
Before the new one: say the Bengali for a riddle, then the Bengali for a trick.
\end{quote}

\subsection*{You'll want to know: \bn{ফাঁদ}}
\textbf{\bn{ফাঁদ}} — \emph{phā̃d} — \textquotedblleft{}a trap\textquotedblright{}.

\begin{quote}
\bn{লোকেরা} \bn{ইঁদুর} \bn{ধরতে} \bn{একটা} \bn{ফাঁদ} \bn{বানাল।}
\end{quote}

\begin{quote}
\emph{lokerā ĩdur dhôrte ekṭā phā̃d bānālô.}
\end{quote}

\begin{quote}
The people made a trap to catch the mice.
\end{quote}

\begin{cousinweb}
\textbf{\bn{ফাঁদ}} (\emph{phā̃d}) comes down from Sanskrit \emph{spāśa}, 'a fetter, a noose' (CDIAL 13813, which lists Bengali \emph{phā̃d}, 'a net'); \textbf{\bn{ফাঁসানো}}, to frame, is its neighbour.
\end{cousinweb}

\begin{grammarlens}[title={What you've built}]
One more word for telling what happened.
\end{grammarlens}

\subsection*{Your turn}
\begin{itemize}
\raggedright
  \item \emph{Say it:} \emph{phā̃d}
  \item \emph{Say it:} \emph{phā̃d}, once more
  \item \emph{Say it:} \emph{lokerā ĩdur dhôrte ekṭā phā̃d bānālô.} — The people made a trap to catch the mice.
  \item \emph{Recall:} say the Bengali for a riddle, then the Bengali for a trick, then say \emph{phā̃d} again
\end{itemize}

\begin{tcolorbox}[breakable,colback=teal!4,colframe=teal!35!black,title={Before you move on}]
What does \bn{ফাঁদ} mean? (\textquotedblleft{}A trap\textquotedblright{}.) Say it once more.
\end{tcolorbox}
\section[\texorpdfstring{\bn{ভোলানো} (bholāno)}{ভোলানো (bholāno)}]{\bn{ভোলানো} (bholāno) — to coax; to beguile}

\label{lesson:BN-C324-bholano}



\begin{quote}
Before the new one: say the Bengali for a trick, then the Bengali for a trap.
\end{quote}

\subsection*{You'll want to know: \bn{ভোলানো}}
\textbf{\bn{ভোলানো}} — \emph{bholāno} — \textquotedblleft{}to coax; to beguile\textquotedblright{}.

\begin{quote}
\bn{শিয়াল} \bn{প্রশংসা} \bn{করে} \bn{কাককে} \bn{ভোলাল।}
\end{quote}

\begin{quote}
\emph{shiyāl prôshôngsā kôre kākke bholālô.}
\end{quote}

\begin{quote}
The jackal flattered the crow and beguiled it.
\end{quote}

Past (he/she …): \bn{সেই} \bn{শিয়াল} \bn{কাককে} \bn{ভোলাল}, \emph{sei shiyāl kākke bholālô}.

\begin{cousinweb}
\textbf{\bn{ভোলানো}} (\emph{bholāno}) is \textbf{\bn{ভোলা}}, to forget — CDIAL 9538, a Middle Indo-Aryan \emph{bhull-}, 'err, forget' — with the -ānō ending: to make someone forget themselves, to coax them.
\end{cousinweb}

\begin{grammarlens}[title={What you've built}]
One more word for telling what happened.
\end{grammarlens}

\subsection*{Your turn}
\begin{itemize}
\raggedright
  \item \emph{Say it:} \emph{bholāno}
  \item \emph{Say it:} \emph{bholāno}, once more
  \item \emph{Say it:} \emph{shiyāl prôshôngsā kôre kākke bholālô.} — The jackal flattered the crow and beguiled it.
  \item \emph{Recall:} say the Bengali for a trick, then the Bengali for a trap, then say \emph{bholāno} again
\end{itemize}

\begin{tcolorbox}[breakable,colback=teal!4,colframe=teal!35!black,title={Before you move on}]
What does \bn{ভোলানো} mean? (\textquotedblleft{}To coax; to beguile\textquotedblright{}.) Say it once more.
\end{tcolorbox}
\section[\texorpdfstring{\bn{পার} \bn{পাওয়া} (pār pāoyā)}{পার পাওয়া (pār pāoyā)}]{\bn{পার} \bn{পাওয়া} (pār pāoyā) — to get away with it; to escape}

\label{lesson:BN-C324-par-paoya}



\begin{quote}
Before the new one: say the Bengali for a trap, then the Bengali for to coax.
\end{quote}

\subsection*{You'll want to know: \bn{পার} \bn{পাওয়া}}
\textbf{\bn{পার} \bn{পাওয়া}} — \emph{pār pāoyā} — \textquotedblleft{}to get away with it; to escape\textquotedblright{}.

\begin{quote}
\bn{চোরটা} \bn{অনেক} \bn{কৌশল} \bn{করল}, \bn{কিন্তু} \bn{পার} \bn{পেল} \bn{না।}
\end{quote}

\begin{quote}
\emph{chorṭā ônek koushôl kôrlô, kintu pār pelô nā.}
\end{quote}

\begin{quote}
The thief tried many tricks, but he didn't get away with it.
\end{quote}

Past (he/she …): \bn{সেই} \bn{চোর} \bn{পার} \bn{পেল}, \emph{sei chor pār pelô}.

\begin{cousinweb}
\textbf{\bn{পার} \bn{পাওয়া}} (\emph{pār pāoyā}) puts \bn{পার}, 'the far bank' — CDIAL 8100, Sanskrit \emph{pāra} — before \textbf{\bn{পাওয়া}}: to reach the far side, to get away. \textbf{\bn{পার} \bn{হওয়া}} is to cross.
\end{cousinweb}

\begin{grammarlens}[title={What you've built}]
One more word for telling what happened.
\end{grammarlens}

\subsection*{Your turn}
\begin{itemize}
\raggedright
  \item \emph{Say it:} \emph{pār pāoyā}
  \item \emph{Say it:} \emph{pār pāoyā}, once more
  \item \emph{Say it:} \emph{chorṭā ônek koushôl kôrlô, kintu pār pelô nā.} — The thief tried many tricks, but he didn't get away with it.
  \item \emph{Recall:} say the Bengali for a trap, then the Bengali for to coax, then say \emph{pār pāoyā} again
\end{itemize}

\begin{tcolorbox}[breakable,colback=teal!4,colframe=teal!35!black,title={Before you move on}]
What does \bn{পার} \bn{পাওয়া} mean? (\textquotedblleft{}To get away with it; to escape\textquotedblright{}.) Say it once more.
\end{tcolorbox}
